\documentclass[10pt,a4paper]{amsart}
\usepackage{amsmath,amssymb,amsthm,graphicx}
\usepackage{verbatim}
\newenvironment{reference}{\comment}{\endcomment}
%\newenvironment{reference}{}{}
\newenvironment{slogan}{\comment}{\endcomment}
\newenvironment{history}{\comment}{\endcomment}

% For commutative diagrams we use Xy-pic
\usepackage[all]{xy}

% We use 2cell for 2-commutative diagrams.
\xyoption{2cell}
% Theorem environments.
%
\theoremstyle{plain}
\newtheorem{theorem}[subsection]{Theorem}
\newtheorem{proposition}[subsection]{Proposition}
\newtheorem{lemma}[subsection]{Lemma}

\theoremstyle{definition}
\newtheorem{definition}[subsection]{Definition}
\newtheorem{example}[subsection]{Example}
\newtheorem{exercise}[subsection]{Exercise}
\newtheorem{situation}[subsection]{Situation}

\theoremstyle{remark}
\newtheorem{remark}[subsection]{Remark}
\newtheorem{remarks}[subsection]{Remarks}

\numberwithin{equation}{subsection}

% Macros
%
\def\lim{\mathop{\mathrm{lim}}\nolimits}
\def\colim{\mathop{\mathrm{colim}}\nolimits}
\def\Spec{\mathop{\mathrm{Spec}}}
\def\Hom{\mathop{\mathrm{Hom}}\nolimits}
\def\Ext{\mathop{\mathrm{Ext}}\nolimits}
\def\SheafHom{\mathop{\mathcal{H}\!\mathit{om}}\nolimits}
\def\SheafExt{\mathop{\mathcal{E}\!\mathit{xt}}\nolimits}
\def\Sch{\mathit{Sch}}
\def\Mor{\mathop{\mathrm{Mor}}\nolimits}
\def\Ob{\mathop{\mathrm{Ob}}\nolimits}
\def\Sh{\mathop{\mathit{Sh}}\nolimits}
\def\NL{\mathop{N\!L}\nolimits}
\def\CH{\mathop{\mathrm{CH}}\nolimits}
\def\proetale{{pro\text{-}\acute{e}tale}}
\def\etale{{\acute{e}tale}}
\def\QCoh{\mathit{QCoh}}
\def\Ker{\mathop{\mathrm{Ker}}}
\def\Im{\mathop{\mathrm{Im}}}
\def\Coker{\mathop{\mathrm{Coker}}}
\def\Coim{\mathop{\mathrm{Coim}}}

% Boxtimes
%
\DeclareMathSymbol{\boxtimes}{\mathbin}{AMSa}{"02}

%
% Macros for moduli stacks/spaces
%
\def\QCohstack{\mathcal{QC}\!\mathit{oh}}
\def\Cohstack{\mathcal{C}\!\mathit{oh}}
\def\Spacesstack{\mathcal{S}\!\mathit{paces}}
\def\Quotfunctor{\mathrm{Quot}}
\def\Hilbfunctor{\mathrm{Hilb}}
\def\Curvesstack{\mathcal{C}\!\mathit{urves}}
\def\Polarizedstack{\mathcal{P}\!\mathit{olarized}}
\def\Complexesstack{\mathcal{C}\!\mathit{omplexes}}
% \Pic is the operator that assigns to X its picard group, usage \Pic(X)
% \Picardstack_{X/B} denotes the Picard stack of X over B
% \Picardfunctor_{X/B} denotes the Picard functor of X over B
\def\Pic{\mathop{\mathrm{Pic}}\nolimits}
\def\Picardstack{\mathcal{P}\!\mathit{ic}}
\def\Picardfunctor{\mathrm{Pic}}
\def\Deformationcategory{\mathcal{D}\!\mathit{ef}}
\usepackage{cleveref}
\crefname{originallemma}{original lemma}{original lemmas}
\Crefname{originallemma}{Original lemma}{Original lemmas}
\makeatletter
\newlabel{lemma-base-change-does-not-hold-pre}{{0F07}{source}{}{sourcepointer-1}{}}
\newlabel{lemma-base-change-does-not-hold-pre@cref}{{[originallemma][0][]0F07}{[1][0][]source}{}{}{}}
\makeatother
\crefname{originallemma}{original lemma}{original lemmas}
\Crefname{originallemma}{Original lemma}{Original lemmas}
\makeatletter
\newlabel{lemma-base-change-q-integral-top}{{0F04}{source}{}{sourcepointer-2}{}}
\newlabel{lemma-base-change-q-integral-top@cref}{{[originallemma][0][]0F04}{[1][0][]source}{}{}{}}
\makeatother
\crefname{originallemma}{original lemma}{original lemmas}
\Crefname{originallemma}{Original lemma}{Original lemmas}
\makeatletter
\newlabel{lemma-base-change-q-integral-bottom}{{0F05}{source}{}{sourcepointer-3}{}}
\newlabel{lemma-base-change-q-integral-bottom@cref}{{[originallemma][0][]0F05}{[1][0][]source}{}{}{}}
\makeatother
\title{A failure of base change is witnessed by an algebraically closed field over a normal affine base}
\author{The Stacks Project Authors}
\date{stacks-project:etale-cohomology.tex:local-snapshot-20261011; GNU FDL 1.2 or later}
\begin{document}
\maketitle
\noindent\textit{Editorial note.} All delivered mathematical text is the author's own, copied from the pinned Stacks Project chapter file etale-cohomology.tex; no statement or proof step was rewritten, reordered or removed.\par\medskip
\noindent\textit{Editorial note.} The proof is the author's complete reduction: the five steps, including the passage to constant cyclic coefficients, are written out and each reduction names the result that licenses it. The author's paragraph at source lines 17319-17323 stating the converse is outside the counted proof span and is not needed for the direction proved here.\par\medskip
\noindent\textit{Editorial note.} Every labelled reference occurring inside the counted spans is carried in assembly.external\_references as a non-load-bearing navigation pointer (same-chapter ones at their definition site, cross-chapter ones at the invocation inside the counted span), so no load-bearing mathematics is replaced by a stub. No bibliography citation occurs inside the counted spans.\par\medskip
\noindent\textit{Editorial note.} The checked chapter list for same-name collisions is the set of stacks chapters staged in this repository: chow, cohomology, divisors, etale-cohomology, flat, more-algebra, more-morphisms, morphisms, spaces-more-morphisms, varieties together with this card's chapter etale-cohomology.\par\medskip
\noindent\textit{Editorial note.} The preamble is the chapter's own native preamble excerpt already used for this source: the Xy-pic setup (native preamble lines 14-18) and the native theorem declarations with the mathematical macros (native preamble lines 154-217).\par\medskip
\noindent\textit{Allowed standard prerequisites (editorial).}\begin{itemize}
\item The base change property BC(f, n, q). Same chapter, Remark remark-base-change-holds, recorded as the context of this card. Supplies the property whose failure is witnessed.
\item Preparation of a diagram where base change fails. Same chapter, Lemma lemma-base-change-does-not-hold-pre. Supplies the starting commutative diagram and the sheaf for which the base change map fails.
\item Base change through an integral top morphism, and through an integral bottom morphism. Same chapter, Lemmas lemma-base-change-q-integral-top and lemma-base-change-q-integral-bottom. Allow the base to be replaced by the normalization in the algebraically closed field without losing the failure.
\end{itemize}\medskip
\ifdefined\hypertarget\hypertarget{sourcepointer-1}{}\fi
\noindent\textit{Source reference.} Original lemma 0F07: sources/stacks-etale-cohomology/pinned-original-source.tex, source lines 17079--17079. Navigation pointer invoked at proof line 17327; the counted span contains the pointer only and the cited result is not reproduced inside it, so the pointer is registered as a non-load-bearing navigation stub in the established form for Stacks material. No bibliography citation is involved.\par\medskip
\ifdefined\hypertarget\hypertarget{sourcepointer-2}{}\fi
\noindent\textit{Source reference.} Original lemma 0F04: sources/stacks-etale-cohomology/pinned-original-source.tex, source lines 16939--16939. Navigation pointer invoked at proof line 17330; the counted span contains the pointer only and the cited result is not reproduced inside it, so the pointer is registered as a non-load-bearing navigation stub in the established form for Stacks material. No bibliography citation is involved.\par\medskip
\ifdefined\hypertarget\hypertarget{sourcepointer-3}{}\fi
\noindent\textit{Source reference.} Original lemma 0F05: sources/stacks-etale-cohomology/pinned-original-source.tex, source lines 16997--16997. Navigation pointer invoked at proof line 17336; the counted span contains the pointer only and the cited result is not reproduced inside it, so the pointer is registered as a non-load-bearing navigation stub in the established form for Stacks material. No bibliography citation is involved.\par\medskip
\begin{remark}
\label{remark-base-change-holds}
Let $f : X \to S$ be a morphism of schemes. Let $n$ be an integer.
We will say $BC(f, n, q_0)$ is true if for every commutative diagram
$$
\xymatrix{
X \ar[d]_f & X' \ar[l] \ar[d]_{f'} & Y \ar[l]^h \ar[d]^e \\
S & S' \ar[l] & T \ar[l]_g
}
$$
with $X' = X \times_S S'$ and $Y = X' \times_{S'} T$ and
$g$ quasi-compact and quasi-separated, and every abelian sheaf
$\mathcal{F}$ on $T_\etale$ annihilated by $n$ the base change map
$$
(f')^{-1}R^qg_*\mathcal{F}
\longrightarrow
R^qh_*e^{-1}\mathcal{F}
$$
is an isomorphism for $q \leq q_0$.
\end{remark}

\begin{lemma}
\label{lemma-base-change-does-not-hold}
With $f : X \to S$ and $n$ as in Remark \ref{remark-base-change-holds}
assume for some $q \geq 1$ we have that
$BC(f, n, q - 1)$ is true, but $BC(f, n, q)$ is not.
Then there exist a commutative diagram
$$
\xymatrix{
X \ar[d]_f & X' \ar[d] \ar[l] & Y \ar[l]^h \ar[d] \\
S & S' \ar[l] & \Spec(K) \ar[l]
}
$$
with both squares cartesian, where
\begin{enumerate}
\item $S'$ is affine, integral, and normal with algebraically
closed function field,
\item $K$ is algebraically closed and $\Spec(K) \to S'$
is dominant (in other words $K$ is an extension of
the function field of $S'$)
\end{enumerate}
and there exists an integer $d | n$
such that $R^qh_*(\mathbf{Z}/d\mathbf{Z})$ is nonzero.
\end{lemma}

\begin{proof}
First choose a diagram and $\mathcal{F}$ as in
Lemma \ref{lemma-base-change-does-not-hold-pre}.
We may and do assume $S'$ is affine (this is obvious, but
see proof of the lemma in case of doubt).
By Lemma \ref{lemma-base-change-q-integral-top}
we may assume $K$ is algebraically closed.
Then $\mathcal{F}$ corresponds to a $\mathbf{Z}/n\mathbf{Z}$-module.
Such a modules is a direct sum of copies of $\mathbf{Z}/d\mathbf{Z}$
for varying $d | n$ hence we may assume $\mathcal{F}$ is
constant with value $\mathbf{Z}/d\mathbf{Z}$.
By Lemma \ref{lemma-base-change-q-integral-bottom}
we may replace $S'$ by the normalization
of $S'$ in $\Spec(K)$ which finishes the proof.
\end{proof}

\end{document}
